\section{Technische und methodische Grundlagen}
\label{sec:grundlagen}

Dieses Kapitel führt die Verfahren ein, die für Architektur,
Versuchsdesign und Interpretation der Pipeline notwendig sind. Die Arbeit
entwickelt keine neuen Segmentierungs-, SfM- oder Splatting-Algorithmen. Ihre
wissenschaftliche Leistung liegt in der konsistenten Verbindung und Prüfung
dieser Verfahren.

\subsection{Promptbare temporale Segmentierung}

Das Segment Anything Model überführt visuelle beziehungsweise sprachliche
Prompts in Objektmasken \parencite{kirillov2023segment,metaSam3}. Für die
Pipeline ist die Promptbarkeit wichtig. Die Zielklasse kann als Begriff wie
\texttt{pipe} angegeben werden, ohne für jede neue lineare Objektart ein
eigenes Modell zu trainieren. Die Videoverarbeitung propagiert die
Objektdefinition über die Sequenz. Damit wird die zeitliche Konsistenz der
Masken zu einem eigenen Qualitätskriterium. Eine gute Einzelbildmaske genügt
nicht, wenn das Objekt in späteren Frames verloren geht.

Der kanonische SAM-Ausgang wird als Hierarchie gespeichert. \texttt{default}
ist die unveränderte Objektmaske, \texttt{middle} entsteht durch eine
$5\times5$-Erosion und \texttt{small} durch zwei Erosionsschritte. Die
Erosion wendet ein quadratisches Strukturelement von $5\times5$ Pixeln an:
Ein Pixel gehört nur dann weiterhin zur Maske, wenn die gesamte
$5\times5$-Nachbarschaft maskiert ist. Dadurch werden unsichere Randpixel
rund zwei Pixel breit entfernt. Der Wert $5\times5$ ist ein pragmatischer
Standard. Er glättet unscharfe Maskenränder und entfernt dünne
Rauschfragmente, ohne das dünne lineare Objekt über eine sinnvolle Bildbreite
hinaus schrumpfen zu lassen. Die konservativeren Stufen entfernen zusätzlich
unsichere Randpixel, können bei sehr dünnen Objekten aber auch relevante
Objektfläche verlieren. Maskenstufe und optionale Dilatation sind deshalb
explizite Parameter. In der aktuellen Projektfassung wird keine Dilatation
angewendet. Historisch war zwischenzeitlich auch die \texttt{default}-Maske
dilatiert. Solche Stände werden als Entwicklungszwischenstände nicht
fortgeführt, sondern durch die verbindliche Festlegung
\texttt{default} = unverändert, \texttt{middle} = einfach erodiert und
\texttt{small} = zweifach erodiert ersetzt (vgl. Abschnitt~\ref{sec:versuchsaufbau}).

\subsection{Structure-from-Motion und Kameramodelle}

Structure-from-Motion schätzt aus korrespondierenden Bildmerkmalen die
Kameraposen und eine sparse 3D-Punktwolke. COLMAP verbindet dafür
SIFT-Merkmalsextraktion, Matching, inkrementelles Mapping und Bundle Adjustment
\parencite{schoenberger2016sfm}. Die Pipeline verwendet unmaskierte Bilder:
Hintergrund- und Ankermerkmale stabilisieren die Orientierung eines glatten,
dünnen Zielobjekts und bleiben für einen späteren GCP-Bezug verfügbar.

Für die Projektarbeit sind drei Kameramodelle relevant. \texttt{PINHOLE}
modelliert getrennte Brennweiten und einen Hauptpunkt, aber keine Verzeichnung. Es setzt damit eine
ideale Bilddomäne voraus. \texttt{SIMPLE\_RADIAL} ergänzt eine gemeinsame
Brennweite um einen radialen Verzeichnungsparameter. \texttt{OPENCV} besitzt
zusätzliche radiale und tangentiale Freiheitsgrade. Das flexiblere Modell ist
nicht grundsätzlich überlegen, da zusätzliche Freiheitsgrade eine
flächendeckende und rauschfreie Merkmalsverteilung voraussetzen. Fehlt diese,
ist das Gleichungssystem schlecht konditioniert – der Optimierer neigt dann
zum Overfitting auf Bildrauschen oder koppelt Verzeichnungsparameter
fälschlicherweise an Kamerabewegungen. Als Produktionsstandard
ist \texttt{OPENCV} gewählt. Die Begründung aus der Matrixauswertung steht in
Abschnitt~\ref{sec:produktionskonfiguration}.

Die Modellbezeichnung ist dabei keine hinterlegte Objektivkalibrierung. Das
Skript \texttt{run\_sfm.sh} übergibt COLMAP lediglich die Modellfamilie und mit
\mbox{\path{ImageReader.single_camera=1}} die Annahme eines gemeinsamen
Kameramodells für alle Frames eines Laufs. Eine externe
\texttt{camera\_params}-Datei oder eine Checkerboard-Kalibrierung wird nicht
eingelesen. COLMAP initialisiert die Parameter intern und schätzt Brennweiten,
Hauptpunkt und Verzeichnungskoeffizienten anschließend gemeinsam mit
Kameraposen und 3D-Punkten aus den Bildkorrespondenzen. Die Werte in
\texttt{sparse\_txt/cameras.txt} sind daher lauf- und auflösungsabhängige
Schätzwerte, keine allgemeinen Standarddaten der Kamera.

Bei verzeichnungsbehafteten Modellen bleibt das geschätzte Originalmodell für
SfM und GCP erhalten. Anschließend erzeugt der COLMAP-
\texttt{image\_undistorter} eine ideale PINHOLE-Szene für STS und die
Meshroute. Bei PINHOLE-Eingaben wird am selben Schnittstellenpfad eine
unveränderte Pass-through-Szene bereitgestellt. Diese Unterscheidung ist
entscheidend: Ein direkter PINHOLE-Lauf auf Rohbildern ist eine Ablation der
Modellannahme und kein Nachweis einer erfolgreichen Entzerrung.

\subsection{3D Gaussian Splatting und Segment-then-Splat}

3D Gaussian Splatting stellt eine Szene als Menge anisotroper Gaussians dar
\parencite{kerbl2023gaussian}, das heißt als Ansammlung kleiner, elliptischer
3D-Verteilungen, deren Form und räumliche Ausrichtung pro Gaussian frei
gelernt werden. Ein Gaussian $i$ besitzt unter anderem einen
Mittelpunkt $\boldsymbol{\mu}_i$, eine räumliche Kovarianz
$\boldsymbol{\Sigma}_i$, Opazität $\alpha_i$ und blickabhängige Farbe. Eine
übliche Zerlegung der Kovarianz lautet

\begin{equation}
	\boldsymbol{\Sigma}_i =
	\mathbf{R}_i\,\mathbf{S}_i\mathbf{S}_i^{\top}\mathbf{R}_i^{\top},
	\label{eq:gaussian-covariance}
\end{equation}

wobei $\mathbf{R}_i$ die Orientierung und $\mathbf{S}_i$ die Skalierung
beschreibt. Beim Rendering werden die Gaussians in die Kamera projiziert,
sortiert und über Alpha-Compositing zu einem Bild zusammengeführt. Die
Repräsentation ist deshalb zugleich geometrische Hypothese und renderbare
Radiance-Field-Darstellung.

Segment-then-Splat ergänzt die Gaussians um objektbezogene Zuordnungen
\parencite{lu2025segment}. Das offizielle Verfahren gewinnt seine
Mehrfachansichtsmasken über eine SAM-basierte Tracking-Pipeline
(AutoSeg-SAM2). Ein SAM-Grid-Prompt segmentiert die erste Ansicht, SAM~2
verfolgt das Objekt anschließend über die Sequenz. Nach der Rekonstruktion
wird jedem Objekt ein CLIP-Embedding zugeordnet, über das die Szene offen
vokabularisch durchsucht werden kann. In der Projektpipeline übernehmen die
extern erzeugten SAM-3.1-Masken beide Rollen. Die promptbare temporale
Segmentierung liefert die Objektmasken direkt, und die CLIP-Embedding-
Zuordnung entfällt, weil für den linearen Einobjektfall ausschließlich die
gewählte Objektmaske als Suchbegriff benötigt wird. STS initialisiert die
Szene aus dem COLMAP-Modell und ordnet jeden Gaussian über die vorgegebenen
Masken einem Objekt zu. Während der Objekt-/Maskenphase wird der Bildverlust
auf die maskierte Objektregion der zugeordneten Maske begrenzt, eine
anschließende All-Object-Phase optimiert alle Objekte gemeinsam. In der
Projektpipeline sind diese Masken die bereits segmentierten Eingabebilder in
den drei Erosionsstufen (\texttt{default}, \texttt{middle}, \texttt{small}),
nicht eine vom STS-System erzeugte Nachsegmentierung. Der Projektstandard
umfasst 7.000 Gesamtiterationen, aufgeteilt in 5.000 Iterationen des
Objektcurriculums und 2.000 Iterationen der All-Object-Phase. Da im vorliegenden
Einobjekt-Szenario keine weiteren Instanzen existieren, bewirkt die finale
Phase keine zusätzliche semantische Entflechtung, sondern dient als konsistente
Trainingsfortführung des isolierten Zielobjekts.

Die Maskenstufen entfalten ihre Wirkung dabei je Komponente getrennt: STS
setzt die bereitgestellten Masken für die Objektzuordnung und die
Objektphasen-Supervision ein. Darüber hinaus ist \texttt{middle} aktuell
allein in der Coverage-Prüfung sowie der Kamera- und Eval-Split-Auswahl
wirksam — Ansichten mit leerer \texttt{middle}-Maske werden vom STS-Training
beziehungsweise von der Evaluation ausgeschlossen. Unabhängig davon belegen
die historischen Zwischenstände in Abbildung~\ref{fig:pa-panel-mesh}, dass
die \texttt{middle}-Stufe die Objektabgrenzung visuell stabilisiert. Der
gezielte Einsatz unterschiedlicher Maskenstufen in den Verlustpfaden der
SuGaR-Vergleichsroute gehört zum maskenbewussten SuGaR-Fork und wird deshalb
zusammen mit diesem Verfahren in Abschnitt~\ref{subsec:meshgewinnung}
beschrieben.

\subsection{Objektselektion und Meshgewinnung}
\label{subsec:meshgewinnung}

Die Ziel-Gaussians werden über die konfigurierte Objekt-ID selektiert. Bei einem
Datensatz mit genau einem linearen Objekt ist die Objekt-ID eine
Formalschnittstelle des übernommenen STS-Schemas. Die eigentliche Auswahl
trifft der nachgelagerte Opazitäts- und Farbfilter. Ein Opazitätsfilter kann
schwach sichtbare Primitive entfernen. Die in der PLY gespeicherte
Logit-Opazität $o_i$ ist der rohe, unbeschränkte Voraktivierungswert, den das
3DGS-Modell lernt. Sie wird erst durch die Sigmoid-Funktion

\begin{equation}
	\alpha_i = \frac{1}{1+\exp(-o_i)}
	\label{eq:opacity}
\end{equation}

in die sichtbare Alpha-Opazität im Bereich $[0,1]$ überführt. Das Filtern über
den Logit erlaubt dabei einen gleichmäßigeren Zugriff auf kleine
Transparenzwerte, als die bereits gesättigte Alpha-Kurve anbieten würde. Nach
Objekt-, Opazitäts- und optionaler Farbauswahl erzeugt die Pipeline eine
getrennte Hochopazitätskopie. Das Setzen auf $\alpha=0{,}999999$ ist eine
Initialisierungspolitik für die geometrische Extraktion, keine physikalische
Aussage über Materialtransparenz: Ein endlicher Wert von exakt 1 ist nach der
Sigmoid-Aktivierung numerisch nicht erreichbar, weil er einen unendlichen
Logit erfordern würde. Der nahezu opake Wert stabilisiert die anschließende
Dichte- und Normalenberechnung beim Surface-Sampling.

SuGaR richtet in seiner vollständigen Route 3D-Gaussians durch
Oberflächenregularisierung als flache Strukturen an einer Oberfläche aus und
rekonstruiert daraus über kamerabasierte Stichproben ein
Poisson-Mesh\footnote{Die Poisson-Oberflächenrekonstruktion legt durch eine
Punktwolke mit orientierten Normalen eine glatte Oberfläche, indem sie ein
Poisson-Gleichungssystem löst, dessen Lösungsfeld die gesuchte
Indikatorfunktion der Oberfläche beschreibt \parencite{kazhdan2006poisson}.}
\parencite{guedon2024sugar}. Im Standardverfahren vergleicht SuGaR stets das
gesamte gerenderte Bild mit dem Trainingsfoto. Da für die Rohrrekonstruktion
jedoch ausschließlich das freigestellte Zielobjekt optimiert werden soll,
begrenzt der im Rahmen dieser Arbeit angepasste lokale SuGaR-Fork den
photometrischen Optimierungsverlust gezielt auf die 2D-Objektmaske:

\begin{equation}
	\mathcal{L}_{\mathrm{RGB}}^{M} =
	\frac{\sum_p M(p)\,\ell_{\mathrm{RGB}}(p)}
	{\sum_p M(p) + \varepsilon}.
	\label{eq:masked-rgb-loss}
\end{equation}

Hierbei ist $M(p) \in \{0, 1\}$ die binäre Objektmaske an Pixelposition $p$
(1 auf dem Rohr, 0 im Hintergrund) und $\ell_{\mathrm{RGB}}(p)$ der kombinierte
Farb- und Strukturfehler ($L_1$-Distanz und DSSIM). Durch die Multiplikation
mit $M(p)$ wird der Fehler im Hintergrund zu null gesetzt, sodass der
Optimierer ausschließlich Gradienten für die Gaussians des Zielobjekts
berechnet und den Hintergrund ignoriert. Die Division durch die Summe aller
Maskenpixel $\sum_p M(p)$ normiert den Verlust auf den Mittelwert pro
Objektpixel. Die winzige Konstante $\varepsilon = 10^{-7}$ schützt vor einer
Division durch null bei temporär leeren Masken.

Beide maskennutzenden Verlustpfade des maskenbewussten SuGaR-Coarse-Trainings
verwenden unterschiedliche Maskenstufen. Die RGB-Supervision nach
Gl.~\ref{eq:masked-rgb-loss} vergleicht das Splat-Rendering jeder
Trainingskamera pixelweise mit dem aufgenommenen Bild, begrenzt auf die
Maskenpixel (maskierter L1-Anteil plus maskierter DSSIM-Anteil mit einem
$11\times11$-Fenster. Das lokale Fenster zieht am Maskenrand begrenzt Kontext
von außen ein). Sie verwendet die unveränderte \texttt{default}-Maske, damit
die Bildsupervision möglichst die volle Objektprojektion abdeckt und
Positionen, Opazitäten und Farben der Objekt-Gaussians vollständig
spezifiziert werden. Eine Farbprojektion auf eine Oberfläche findet dabei
nicht statt. Die Farbübertragung aus den Ansichten erfolgt erst in der
nachgelagerten UV-Texturphase des Forks.

Die Depth-Normal-Consistency (DN) vergleicht dagegen die aus der gerenderten
Gaussian-Tiefe abgeleiteten Normalen mit den gerenderten Normalen und wird
bewusst auf der erodierten \texttt{middle}-Stufe ohne Dilatation angewandt.
Der Grund ist die Silhouettenstabilität. Tiefe und daraus abgeleitete Normalen
sind direkt an Objektkanten unzuverlässig und genau dort kann eine
\texttt{default}-Maske über die Objektsilhouette hinausragen. Ein solcher
Maskenüberhang würde die Objektnormalen am Silhouettenrand in Richtung der
Hintergrundtiefe ziehen und dadurch eine ausgestülpte oder abgeflachte
Oberfläche erzeugen. \texttt{middle} entfernt diese unsichere Randzone von
rund zwei Pixeln, sodass nur Pixel supervisiert werden, die sicher auf dem
Objekt liegen.

Diese Aufteilung stammt aus der Entwicklung des Forks und ist dort Teil des
Datenvertrags zwischen Maskenhierarchie und Trainingsplan. Da die
DN-/SDF-Terme erst oberhalb des Zählerstands 9.000 aktiviert werden und alle
ausgewerteten Läufe — einschließlich sämtlicher \texttt{sugar\_coarse}-Arme
der Vierfeld-Ablation und der Folgeprüfung — mit Zielzählerstand 9.000 oder
weniger trainiert wurden, trat der DN-Term in keinem ausgewerteten Lauf
tatsächlich in Wirkung. Die \texttt{middle}-Zuordnung bleibt dennoch
dokumentiert, weil sie den Fork-Vertrag definiert. Zählermechanik und
Nachweislage der Läufe sind in Abschnitt~\ref{subsec:sugar-folgepruefung}
erläutert. Standardläufe mit Zieliteration 9.000 optimieren die Geometrie
daher primär über den maskierten photometrischen Farbverlust.


Auf dieser Basis definiert das Projekt eine direkte Meshroute, die ohne
Coarse-Training auskommt: Sie überspringt die zusätzliche SuGaR-Coarse-
Optimierung, das Gaussian-bound Refinement und das UV-Baking. Sie verwendet
Original-STS-Gaussians mit projizierter Diamond-Mesh-Tiefe\footnote{Die
Diamond-Mesh-Tiefe wird aus der Projektion der SuGaR-``Diamond''-Primitive
gewonnen: Diese flachen, an den Gaussians ausgerichteten Vierecksflächen
werden in einen Z-Buffer projiziert, aus dem je Pixel eine Tiefe abgelesen
wird; das Surface-Sampling wandelt diese Tiefe anschließend in 3D-Punkte. Die
Alternative ist die Tiefe direkt aus dem Gaussian-Rasterizer.}, Surface-
Sampling, Poisson-Rekonstruktion, Dichtebereinigung, Decimation und optionale
Vertexprojektion. Die Dichtebereinigung entfernt Cluster mit zu wenigen
Stützpunkten, die Decimation reduziert die Dreieckszahl auf das Ziel-Vertice-
Budget, und die optionale Vertexprojektion verschiebt die verbliebenen
Vertices auf die nächstgelegenen Oberflächenstichproben beziehungsweise deren
Tiefenpunkte. Diese Konfiguration entspricht genau der Zelle A der
kontrollierten Vierfeld-Ablation (Abschnitt~\ref{subsec:vierfeld-ablation});
nach dieser Zelle wird sie im Folgenden durchgängig als \emph{Route A}
bezeichnet, während die Coarse-Route entsprechend dem Modusnamen des
Ausführungscodes \texttt{sugar\_coarse} heißt; Route A selbst trägt dort den
Modusnamen \texttt{original\_gs}. Die vollständige SuGaR-Coarse-Route bleibt als
wissenschaftlicher Vergleich erhalten. Damit werden zwei Fragen getrennt:
Wie wird aus einer Gaussian-Wolke ein Mesh gewonnen, und wie verändert eine
zusätzliche Coarse-Optimierung zuvor die Gaussian-Geometrie? SuGaR ist
konzeptionell so ausgelegt, dass Meshes über seine volle Route (Coarse-
Optimierung mit anschließender Oberflächenextraktion) entstehen. Die in der
Projektarbeit verwendete Route A ist deshalb die Abweichung vom
Standardweg: Sie extrahiert direkt aus der unveränderten Original-STS-
Gaussian-Wolke, ohne die Coarse-Optimierung. Diese schlankere Variante wurde
erst im Projektverlauf ergänzend erprobt und hat anschließend die
Produktionsrolle übernommen; die SuGaR-Coarse-Route wird als
Vergleichsarm fortgeführt. Beide Routen bleiben als getrennte und
kombinierbare Bausteine denkbar: Eine spätere Variante könnte auf denselben
Eingaben zusätzlich zur segmentierten Objektrekonstruktion auch eine
vollständige Szene (Gaussians und Mesh) erzeugen, ohne den Objektpfad zu
verändern.

\subsection{Centerline und B-Spline}

DGtal voxeliert das Mesh auf ein diskretes Gitter (Voxelgröße 0,1 in den
lokalen Skaleneinheiten des unskalierten Modells; ohne reale Maßreferenz ist
dies eine relative, keine metrische Größe) und führt eine
topologieerhaltende Skelettierung aus
\parencite{dgtalProject}. Auf verrauschten, dünnwandigen Oberflächen können
dabei Medialflächen (flächige Skelettbereiche, für eine Linientrasse nicht
direkt verwertbar; vgl. die Fußnote in Kapitel~\ref{sec:datengrundlage}) und
zahlreiche kurze Äste entstehen. Der produktive
\texttt{single}-Modus verwendet deshalb den Durchmesserpfad der größten
Skelettkomponente. Dies liefert eine robuste Einzeltrasse, bildet aber keine
verzweigten Netze ab.

Die diskrete Centerline wird durch einen geklemmten uniformen B-Spline des
Grades $p=10$ geglättet. Die Kurve ist

\begin{equation}
	\mathbf{C}(u)=\sum_{i=0}^{n}N_{i,p}(u)\,\mathbf{P}_i,
	\label{eq:bspline}
\end{equation}

mit Kontrollpunkten $\mathbf{P}_i$ und B-Spline-Basisfunktionen $N_{i,p}$.
Der relativ hohe Grad $p=10$ glättet die diskrete Centerline ausreichend,
damit einzelne Ausreißer nicht zu zackigen Kurvensprüngen führen; er ist ein
algorithmischer Glättungsparameter, der im Pipeline-Skript manuell einstellbar
ist, und begründet keine metrische Genauigkeit. Die Anzahl der ausgegebenen
Samples verändert nur die Diskretisierung der Kurve.

\subsection{Objektmaskierte Bildmetriken und Splat-Checkpoints}

Die Bildmetriken werden nicht auf dem Dreiecksnetz (Mesh), sondern auf den
renderbaren 3D-Gaussian-Splatting-Repräsentationen ausgewertet. Dafür werden
synthetische Ansichten für einen festen, beim Training ausgeschlossenen
Test-Kamerasplit (\path{eval_frames.txt}) gerendert und pixelweise mit den
echten Referenzbildern verglichen.

\subsubsection*{Splat-Serialisierung: PLY- versus PT-Dateien}

In der Pipeline treten zwei verschiedene Speicherformate für Gaussian-Checkpoints
auf:
\begin{itemize}
  \item \textbf{PLY-Dateien (\path{.ply}):} Das standardisierte Polygon File
        Format wird von Vanilla 3DGS und STS (\path{point_cloud.ply}) verwendet. Es
        speichert die Gaussian-Eigenschaften als Punktattribute (Positionen $x,y,z$,
        Normalen und Rotationsquaternionen, Skalierungsfaktoren, Logit-Opazität und
        Spherical-Harmonics-Koeffizienten) und ist in externen Visualisierern (z.\,B.
        CloudCompare oder SuperSplat) direkt darstellbar.
  \item \textbf{PyTorch-Checkpoints (\path{.pt}):} SuGaR speichert
        Trainingszwischenstände der Coarse-Stufe standardmäßig als serialisiertes
        PyTorch-\texttt{state\_dict} (z.\,B. \path{9000.pt}), da SuGaR als
        \texttt{torch.nn.Module} implementiert ist.
\end{itemize}

Beide Formate beschreiben bei identischem Trainingsstand dieselben
3D-Gaussian-Parameter ($\boldsymbol{\mu}_i, \mathbf{S}_i, \mathbf{R}_i,
\alpha_i$ und SH-Farbwerte); bei vollständiger Attributübernahme in gleicher
Präzision ist ein Laden aus einer \path{.pt}-Datei in das SuGaR-Modell
inhaltlich gleichwertig zu einem Export in eine standardkonforme
\path{.ply}-Datei. Unterschiede in den resultierenden Bildmetriken entstehen
daher aus dem unterschiedlichen Trainingszustand der Gaussians (z.\,B. 7.000
Iterationen STS gegenüber 9.000 Iterationen SuGaR-Coarse-Optimierung), nicht
aus dem Dateiformat.

\subsubsection*{Berechnung der Metriken}

Die Bildmetriken werden ausschließlich objektmaskiert aggregiert. Für $N_M$ gültige
Maskenpixel und normierte RGB-Werte wird der kanalweise gemittelte maskierte
mittlere quadratische Fehler als

\begin{equation}
	\operatorname{MSE}_M=
	\frac{1}{3N_M}\sum_{p:M(p)=1}\lVert I(p)-\hat I(p)\rVert_2^2
	\label{eq:masked-mse}
\end{equation}

berechnet. Dabei ist $I(p)$ das gerenderte Bild der Gaussian-Repräsentation an
Pixelposition $p$ und $\hat I(p)$ das zugehörige reale Referenzbild (Ground
Truth) aus dem festen Test-Kamerasplit; summiert wird nur über Pixel mit
$M(p)=1$, also innerhalb der Objektmaske. Daraus folgt bei einem Maximalwert
von eins

\begin{equation}
	\operatorname{PSNR}_M=10\log_{10}\!\left(\frac{1}{\operatorname{MSE}_M}\right).
	\label{eq:masked-psnr}
\end{equation}

Ein höherer PSNR-Wert bedeutet eine geringere pixelweise Abweichung. SSIM
bewertet zusätzlich lokale Helligkeit, Kontrast und Struktur; höhere Werte
sind besser \parencite{wang2004ssim}. Die lokale Fenstergröße von
$11\times11$ Pixeln ist dabei der Literaturstandard nach Wang et al. (2004)
\parencite{wang2004ssim}. Im aktuellen Evaluator wird zunächst eine
vollständige SSIM-Karte berechnet und anschließend nur an gültigen
Maskenpixeln aggregiert; der lokale SSIM-Fensterkontext kann daher weiterhin
Randpixel außerhalb der Maske enthalten. Bei linearen Objekten fällt dieser
Randeffekt stärker ins Gewicht, je kleiner das Objekt im Pixelraster ist:
Rund fünf Randpixel bedeuten bei einem etwa 40 Pixel breiten Rohr (720p)
nur einen Randanteil von etwa 12 Prozent, bei einem 20 Pixel breiten Rohr
(Low-Auflösung) jedoch bereits etwa 25 Prozent. LPIPS vergleicht Merkmale eines
vortrainierten Netzes; niedrigere Werte bedeuten eine größere perzeptuelle
Ähnlichkeit \parencite{zhang2018lpips}. PSNR und SSIM bewerten stets nur das
Signal im jeweiligen Pixelgitter der Renderauflösung – sie liefern kein
absolutes Weltkoordinatenmaß. Full-frame-Metriken sind
für den object-only Splat deaktiviert, weil sie überwiegend den nicht
rekonstruierten Hintergrund bewerten würden.

PSNR, SSIM und LPIPS sind Ansichtsmetriken. Sie beweisen weder eine korrekte
Meshoberfläche noch eine genaue Centerline. Ebenso garantiert eine vollständige
Maske nicht, dass Kamerageometrie und Surface-Sampling metrisch korrekt sind.
Geometrische RMSE- und Hausdorff-Auswertungen gegen eine unabhängige
GNSS-Referenz bleiben Gegenstand der Bachelorarbeit.
